\subsection{Уголовный кодекс Российской Федерации}

Уголовный кодекс Российской Федерации \cite{UKRF} осуществляет уголовно-правовую защиту информации, информационных систем и отношений, возникающих в сфере их использования, устанавливая ответственность за совершение преступлений в данной области. В контексте обеспечения информационной безопасности особое значение имеют положения главы 28 УК РФ, посвящённой преступлениям в сфере компьютерной информации.

Статьи 272–274, а также 272.1 УК РФ регламентируют ответственность за нарушения, связанные с несанкционированным доступом, неправомерным использованием компьютерной информации и эксплуатацией средств её обработки.

Статья 272.1 УК РФ предусматривает уголовную ответственность за незаконное использование, передачу, сбор или хранение компьютерной информации, содержащей персональные данные, полученные незаконным путём, а также за создание и функционирование ресурсов, предназначенных для их незаконного распространения. Данная норма направлена на защиту персональных данных и пресечение их нелегального оборота, включая случаи утечек, торговли базами данных и иных форм неправомерного обращения с такой информацией.

Статья 273 УК РФ устанавливает ответственность за создание, использование и распространение вредоносных программ для ЭВМ, предназначенных для несанкционированного уничтожения, блокирования, модификации или копирования информации, а также для нейтрализации средств защиты информации.

Статья 274 УК РФ регламентирует ответственность за нарушение правил эксплуатации средств хранения, обработки или передачи компьютерной информации, а также информационно-телекоммуникационных сетей, если это повлекло уничтожение, блокирование или модификацию информации и причинило существенный вред.